% ==============================================================
% Kapitola: Záver
% ==============================================================

\section{Záver}
\label{sec:conclusion}

Táto práca sa zaoberala návrhom, implementáciou a experimentálnym hodnotením bezpečnostných mechanizmov pre systém kontroly fyzického prístupu postavený na bezkontaktných identifikátoroch RFID/NFC. Výsledkom je funkčný prototyp a jeho systematická analýza, ktorá identifikuje konkrétne prínosy jednotlivých bezpečnostných vrstiev.

\subsection{Súhrn realizovaných prác}

V rámci práce boli realizované nasledujúce kroky:

\begin{enumerate}[leftmargin=*,itemsep=3pt]
  \item \textbf{Teoretická analýza.} Rozobrané boli bezpečnostné pojmy (CIA triáda, AAA model, RBAC), kryptografické primitíva (AEAD, HMAC, HKDF), návrh bezpečných protokolov (fail-closed, kanonizácia, anti-replay) a špecifiká RFID/NFC technológií vrátane známych útokov na MIFARE Classic\cite{GarciaMIFARE2008}. Analýza bola orientovaná na prax --- každý teoretický koncept je priamo naviazaný na konkrétny mechanizmus v prototype.

  \item \textbf{Návrh a implementácia prototypu.} Prototyp pozostáva z hardvérovej čítačky (ESP32 + RC522\cite{Espressif2024}), serverovej aplikácie v jazyku C++20 a lokálneho úložiska SQLite. Systém implementuje viacvrstvovú ochranu:
  \begin{itemize}[itemsep=1pt]
    \item HMAC-SHA-256 autentifikácia hardvérových požiadaviek,
    \item interný AEAD frame (ChaCha20-Poly1305\cite{RFC8439} / XChaCha20-Poly1305\cite{XChaChaDraft}) s AAD väzbou hlavičky,
    \item HKDF derivácia per-reader kľúčov\cite{RFC5869} s verziovaním a domain separation,
    \item tri stratégie generovania nonce (náhodný, deterministický HMAC, deterministický s runtime verifikáciou),
    \item \texttt{ProtocolAnomalyDetector} so schopnosťou karantény zariadení,
    \item sliding window anti-replay mechanizmus,
    \item tamper-evident auditný log s HMAC hash chain\cite{SchneierKelsey1999} a anchor ochranou,
    \item RBAC model s pepper-chránenými identifikátormi kariet.
  \end{itemize}
  Kódová báza je modulárna (9~modulov), písaná v štandarde C++20 s výhradným použitím knižnice libsodium pre kryptografické operácie\cite{LibsodiumXChaChaDoc}.

  \item \textbf{Verifikácia kvality kódu.} Pred experimentmi bola vykonaná statická analýza (\texttt{clang-tidy}\cite{ClangTidy}: 0~bezpečnostných nálezov) a dynamická analýza (ASan\cite{AddressSanitizer}, UBSan, LSan: 0~nálezov za behu). Testovacia sada\cite{GoogleTest} obsahuje 44~testovacích prípadov pokrývajúcich kryptografiu, replay window, auditný log a rozhodovaciu logiku.

  \item \textbf{Experimentálna analýza.} Navrhnutý a implementovaný bol C++ experimentálny harness, ktorý zdieľa produkčný kód a umožňuje reprodukovateľné testy. Šesť konfiguračných profilov (2~algoritmy $\times$ 3~nonce stratégie) bolo testovaných v šiestich scenároch (replay, AAD tampering, cross-reader key isolation, seq rollback, tag forgery resistance, throughput benchmark). Každý scenár bežal 5-krát, celkový objem dát presahuje 2{,}8~milióna meraní.

  \item \textbf{Analýza a vizualizácia výsledkov.} Python pipeline spracúva CSV výstupy do siedmich grafov a troch LaTeX tabuliek s~chybovými pásmami a~intervalmi spoľahlivosti.
\end{enumerate}

\subsection{Hlavné výsledky a závery}

Experimentálna analýza potvrdila šesť hlavných tézisov:

\begin{table}[htbp]
\renewcommand{\arraystretch}{1.3}
\centering
\caption{Súhrn hlavných experimentálnych zistení}
\label{tab:findings}
\small
\begin{tabular}{c p{10.5cm}}
\toprule
\textbf{Tézis} & \textbf{Zistenie} \\
\midrule
T1 & R2 ako jediný vynucuje nonce politiku (S3a: R0/R1 100\% úspešnosť vs.\ R2 0\%) --- bezpečnostná poistka pre budúce škálovanie. HKDF per-reader derivácia zabraňuje cross-reader útokom (S3b: 0\% úspešnosť naprieč všetkými profilmi). \\
T2 & \texttt{ProtocolAnomalyDetector} (profily R2) dosiahol 0\% úspešnosť naprieč všetkými testovanými scenármi (S1--S5). V S3a a S4 bol jediným mechanizmom, ktorý scenár úplne zastavil; v S3b a S5 pridáva operačnú eskaláciu (karanténu). \\
T3 & Replay window a anomálna detekcia sú komplementárne. S4 odhaľuje štruktúrnu medzeru replay window (56 z 500 frame prejde), ktorú detektor uzatvára cez rollback detekciu. \\
T4 & AAD binding poskytuje ochranu nezávisle od režimu nonce --- S2 ukazuje 0\% úspešnosť manipulácie hlavičky naprieč všetkými profilmi. \\
T5 & Karanténa zariadení mení reakciu systému z reaktívnej (odmietnutie frame) na proaktívnu (izolácia zariadenia) --- realizácia defense-in-depth\cite{DefenseInDepth}. \\
T6 & Overhead detektora (${\sim}19\%$ pokles priepustnosti, ${\sim}35\,\mu$s na frame) je rádovo nepodstatný pre reálnu prevádzku prístupového systému ($<$10 udalostí/s), hoci nie je zanedbateľný. \\
\bottomrule
\end{tabular}
\end{table}

\noindent Kľúčovým zistením je, že profil R2 (deterministický nonce s aktívnym detektorom) v kombinácii s~existujúcimi vrstvami (AEAD + AAD + HKDF + anti-replay + audit) poskytuje \textbf{materiálne vyššiu odolnosť} voči všetkým testovaným triedam útokov pri prijateľnom výkonnostnom overheade. Zároveň experiment potvrdzuje, že \textbf{voľba algoritmu} (A1 vs.\ A2) nemá merateľný vplyv na bezpečnostné metriky --- oba algoritmy sú ekvivalentné v odolnosti a porovnateľné vo výkone.

Všetkých 5 behov (seedy 42--46) poskytlo zhodné výsledky bezpečnostných metrík, čo naznačuje, že pozorované rozdiely vychádzajú z architektonických vlastností systému, nie z konkrétnych hodnôt nonce.

\subsection{Splnenie cieľov práce}

\begin{table}[htbp]
\renewcommand{\arraystretch}{1.3}
\centering
\caption{Mapovanie cieľov práce na realizované výstupy}
\label{tab:goals-mapping}
\begin{tabular}{c p{5.8cm} p{5.8cm}}
\toprule
\textbf{Cieľ} & \textbf{Definícia} & \textbf{Realizácia} \\
\midrule
C1 & Analýza bezpečnostných hrozieb & Kapitoly 3--8: CIA triáda, STRIDE\cite{STRIDE}, RFID/NFC hrozby, kryptografické základy \\
C2 & Návrh bezpečnostného protokolu & Kapitoly 9--10: HMAC autentifikácia, AEAD frame, anti-replay, audit \\
C3 & Implementácia prototypu & Kapitola 21: 9 C++ modulov, ESP32 firmware, SQLite úložisko \\
C4 & Verifikácia kvality kódu & Kapitola 22: clang-tidy, sanitizers, 44 testov \\
C5 & Experimentálne hodnotenie & Kapitola 23: 6 profilov $\times$ 6 scenárov $\times$ 5 behov \\
\bottomrule
\end{tabular}
\end{table}

\subsection{Otvorené otázky a obmedzenia}

Napriek pozitívnym výsledkom je potrebné identifikovať obmedzenia, ktoré ovplyvňujú interpretáciu a prenášateľnosť záverov:

\begin{enumerate}[leftmargin=*,itemsep=3pt]
  \item \textbf{UID nie je kryptografický dôkaz identity.} Karty typu MIFARE Classic\cite{GarciaMIFARE2008} alebo iné karty bez vzájomnej autentifikácie umožňujú emuláciu UID. Prototyp sa spolieha na to, že bezpečnosť je zabezpečená na úrovni protokolu (HMAC + AEAD), nie na úrovni karty. Nasadenie kryptograficky silných kariet (napr. MIFARE DESFire) by eliminovalo túto závislosť.

  \item \textbf{Absencia TLS.} Komunikácia medzi čítačkou a serverom nie je šifrovaná na transportnej vrstve. V lokálnej sieti je HMAC autentifikácia dostatočná na zabránenie podvrhnutiu, no v sieťach s nedôveryhodnou infraštruktúrou by bolo nutné doplniť TLS\cite{RFC8446} alebo mTLS.

  \item \textbf{Jednostranná autentifikácia.} Server autentifikuje čítačku (cez HMAC), ale čítačka neautentifikuje server. Man-in-the-middle útočník by teoreticky mohol zachytávať odpovede. Pre produkčné nasadenie by bolo vhodné zvážiť vzájomnú autentifikáciu\cite{RFC9325}.

  \item \textbf{Lokálna SQLite databáza.} V aktuálnom riešení je databáza uložená na rovnakom stroji ako server. Útočník s fyzickým prístupom k serveru môže modifikovať databázu aj audit anchor. Produkčné nasadenie by vyžadovalo vzdialený anchor alebo HSM\cite{NIST80057}.

  \item \textbf{Rozsah testovaných scenárov.} Šesť scenárov pokrýva hlavné triedy útokov, no nepokrýva napríklad timing útoky, side-channel analýzu, alebo útoky na firmware čítačky. Rozšírenie sady scenárov by zvýšilo dôveru v kompletnosť hodnotenia.

  \item \textbf{Jednoprocesový server.} Aktuálna implementácia beží ako jeden proces s jedným vláknom pre spracovanie požiadaviek. Škálovanie na viacero čítačiek v reálnom nasadení by vyžadovalo asynchronné spracovanie alebo viacprocesový model.
\end{enumerate}

\subsection{Možnosti rozšírenia}

Na základe identifikovaných obmedzení a skúseností z implementácie je možné navrhnúť niekoľko smerov budúceho vývoja:

\begin{description}[leftmargin=*,itemsep=3pt]
  \item[Integrácia TLS/mTLS:] Doplnenie šifrovanej transportnej vrstvy s certifikátovou autentifikáciou zariadení\cite{RFC8446,RFC9325}. ESP32 podporuje TLS cez mbedTLS, čo umožňuje implementáciu bez zmeny hardvéru.

  \item[Podpora kryptografických kariet:] Prechod z jednoduchého čítania UID na vzájomnú autentifikáciu s kartou (napr. MIFARE DESFire EV3). Tým by sa eliminovala závislosť na serverovej kompenzácii slabej identity karty.

  \item[Fuzz testing:] Systematické testovanie vstupných parserov a kryptografických rozhraní pomocou AFL++ alebo libFuzzer s~cieľom nájsť edge cases, ktoré jednotkové testy nepokrývajú.

  \item[Distribuovaná architektúra:] Rozšírenie na viacero serverov s centralizovanou správou kľúčov a federovaným auditom. Anchor by bol uložený na nezávislom serveri alebo v~blockchain-like štruktúre.

  \item[Rotácia pepperu s online migráciou:] Implementácia bezprestojovej migrácie pepper verzie, kde sa existujúce karty postupne re-hashujú pri prvom použití pod novou verziou.

  \item[Adaptívna karanténa:] Rozšírenie \texttt{ProtocolAnomalyDetector} o časovo obmedzenú karanténu s~eskaláciou (napr. 1. incident: 5~minút; 2. incident: 1~hodina; 3. incident: trvalá karanténa). Aktuálna implementácia používa trvalú karanténu, čo je bezpečné, ale menej flexibilné.

  \item[Offline režim čítačky:] Podpora scenárov, kde čítačka dočasne stratí spojenie so serverom. Toto vyžaduje lokálnu cache posledných autorizácií a následnú synchronizáciu, čo prináša nové bezpečnostné výzvy (stale cache, conflict resolution).

  \item[Monitoring a alerting:] Integrácia s existujúcimi monitorovacími systémami (Prometheus, Grafana) pre vizualizáciu prevádzkových metrík a automatické upozornenia pri detekcii anomálií\cite{NIST80094}.
\end{description}

\subsection{Záverečné zhodnotenie}

Práca demonštruje, že aj v prototypovom prostredí je možné implementovať viacvrstvový bezpečnostný systém, ktorý poskytuje merateľnú ochranu voči realistickým útokovým scenárom. Kľúčovým prínosom nie je samotná implementácia jednotlivých mechanizmov --- tie sú v literatúre dobre popísané --- ale ich \emph{systematická integrácia} a \emph{experimentálne overenie} interakcií medzi nimi.

Výsledky jasne ukazujú, že bezpečnosť nie je binárna vlastnosť (\uv{zabezpečené vs.\ nezabezpečené}), ale spektrum závislé od konkrétnej kombinácie mechanizmov. Profil R2 s~aktívnym anomálnym detektorom nie je \uv{lepší} v~abstraktnom zmysle --- je lepší v~presne definovaných scenároch (S3, S4), kde iné mechanizmy zlyhávajú. Iné mechanizmy (replay window, AAD binding, AEAD) sú zasa dostatočné v~scenároch, kde detektor nie je potrebný (S1, S2, S5).

Tento diferencovaný pohľad --- nie \uv{zapni všetko}, ale \uv{pochop, čo chráni proti čomu} --- je hlavným záverom experimentálnej časti a podľa záverov analýzy aj najdôležitejším prínosom tejto práce pre prax návrhu bezpečných systémov kontroly prístupu.
